% PreTextuais/Resumo.tex
\begin{resumo}
A gestão de frotas no setor público historicamente opera sob rotinas reativas emergenciais devido à fragmentação de dados entre cadastros patrimoniais, ordens de serviço (OSs) e abastecimentos. Esta dissertação propõe e avalia uma arquitetura computacional de analítica prescritiva orientada por Inteligência Artificial Explicável (XAI), Teoria dos Grafos e Otimização Inteira Mista (MILP). A base empírica abrange 77.867 ordens de serviço transacionadas entre 2020 e 2026 no Poder Executivo de Santa Catarina, totalizando R\$ 138,76 milhões em dispêndio financeiro aprovado e liquidado (após a retenção de R\$ 60,05 milhões em glosas técnicas preventivas pela fiscalização contratual), vinculadas a 34.280 veículos e 32,84 milhões de litros de combustível. Avaliaram-se regressores e classificadores não lineares sob validação cruzada temporal. No Cenário Canônico S0 (2.000 veículos monitorados por 36 períodos mensais), o modelo MILP resolvido pelo solver HiGHS com gap de otimalidade de 0,0000\% alcançou uma economia simulada de R\$ 9.302.147,26 (17,51\% frente à política inercial), com 100\% de conformidade com manutenções obrigatórias. A fusão multimodal (H1) apresentou significância estatística, mas ganho marginal de 0,41\% no erro médio absoluto, sendo classificada como não sustentada frente ao patamar de relevância prática pré-fixado de 15,0\%. A avaliação de explicabilidade (H4) revelou concordância Top-5 de Jaccard de 39,29\% entre SHAP e LIME, não sustentando a intercambialidade direta dos métodos e fundamentando a necessidade de protocolos de triangulação analítica para auditoria governamental. Por fim, identificaram-se 313 veículos em condição antieconômica para subsidiar decisões de alienação e leilão. Conclui-se que a analítica prescritiva sob restrições contrafactuais auditadas viabiliza ganhos fiscais mensuráveis e transparência decisória para a governança pública.

\vspace{\onelineskip}
\noindent
\textbf{Palavras-chave:} analítica prescritiva; gestão de frotas públicas; teoria dos grafos; inteligência artificial explicável; programação linear inteira mista.
\end{resumo}
